\section{Gap decay on both sides of the Liouville threshold}\label{sec:intro}
\paragraph{NEW.}
For a supercritical almost Mathieu operator, positive Liouville exponent changes the gap-decay picture even below the threshold $2\beta=\log\lambda$. The same line is the spectral transition at the completely resonant phases, $2\theta\in\alpha\Z+\Z$: Liu proved Anderson localization there for $\log\lambda>2\beta$~\cite{Liu2022} and, by a doubled Gordon threshold, purely singular continuous spectrum for $1<\lambda<e^{2\beta}$~\cite{Liu2026}. We prove an upper bound on the normalized decay rate, $\rho\le1-2\beta/\log\lambda$, below that threshold. Together with the polynomial lower subsequences established in version~1.0.0 of this paper~\cite{PaperDv1}, this gives quantitative gap-width results on both sides. The equality line and the matching lower bound below it remain open.

\paragraph{KNOWN (normalization and gap labels).}
We use hopping one:
\begin{equation}\label{eq:H}
(H_{\lambda,\alpha,\theta}u)(j)=u(j+1)+u(j-1)
 +2\lambda\cos(2\pi(j\alpha+\theta))u(j).
\end{equation}
Throughout, $\lambda>1$, $L=\log\lambda$, and $\alpha$ is irrational. Its convergents are $p_n/q_n$, and
\begin{equation}\label{eq:arithmetic}
\delta_n=q_n\alpha-p_n,\quad b_n=-\frac{\log|\delta_n|}{q_n},
\quad \beta=\limsup_n\frac{\log q_{n+1}}{q_n}=\limsup_n b_n.
\end{equation}
The common spectrum is $\Sig$. The integrated density of states $N$ labels a bounded gap $G_m$ by $N|_{G_m}=\{m\alpha\}$, $m\ne0$; see~\cite{AvronSimon,AYZdry}. We normalize the gap exponent as
\begin{equation}\label{eq:rho}
\rho(\lambda,\alpha)=\liminf_{|m|\to\infty}
 \frac{-\log|G_m|}{L|m|},\qquad x=\frac{2\beta}{L}.
\end{equation}
A closed gap contributes $+\infty$. Bounded spectral diameter gives $\rho\ge0$.

\paragraph{KNOWN (rational scales).}
Let $P_n$ be the monic Chambers polynomial and $E_n$ the top of the rational phase union. Set
\begin{equation}\label{eq:scales}
\Lambda_n=\lambda^{q_n},\qquad
\hbar_n=2\pi q_n\Lambda_n^{1/2}|\delta_n|,\qquad
\kappa_n=\hbar_n^2.
\end{equation}
The derivative $P_n'(E_n)$ is positive; see Lemma~\ref{lem:convex}.
\begin{figure}[ht]
\centering\small
\fbox{\begin{minipage}{.93\linewidth}
\centering
Rational top band\quad$\longrightarrow$\quad Translated wells\quad$\longrightarrow$\quad Common ladder\\[4pt]
$P(E_*)=2\Lambda+2$\qquad
one well per $|\delta|^{-1}$ sites\qquad
spacing of order $\hbar/P_*'$\\[7pt]
$D$ levels above $E$ in every well\quad$\Longrightarrow$\quad
$N(E)=1-D|\delta|$\\[3pt]
$\Longrightarrow$\quad a gap with label $m=-Dq\sgn\delta$.
\end{minipage}}
\caption{The single-scale mechanism, with the scale index omitted and $P_*'=P'(E_*)$. The common well count is proved before the infinite-volume IDS limit.}\label{fig:mechanism}
\end{figure}


\paragraph{NEW.}
The following is the new result of version~2.0.0.
\begin{maintheorem}[Fixed-label-multiple gaps below the threshold]\label{thm:below}
Suppose $0<\beta<L/2$. Fix an integer $D\ge1$. There is $c_D>0$, depending only on $D$, such that, on any subsequence of convergents with $b_n\to\beta$, for all sufficiently large $n$,
\begin{equation}\label{eq:belowwidth}
 |G_{-Dq_n\sgn\delta_n}|
 \ge \frac{c_D\kappa_n^{1-D}}{P_n'(E_n)}.
\end{equation}
The threshold in $n$ may depend on $\lambda,\alpha,\beta,D$. In particular, for every $\epsilon>0$, eventually on that subsequence,
\begin{equation}\label{eq:belowexponential}
 |G_{-Dq_n\sgn\delta_n}|
 \ge \exp\!\left(-\bigl[(D-1)(L-2b_n)+L+\epsilon\bigr]q_n\right).
\end{equation}
Consequently $\rho(\lambda,\alpha)\le1-x$ for $0<x<1$.
\end{maintheorem}
No uniformity in growing $D$ is asserted here. The order of limits is first $n\to\infty$ at fixed $D$, then $D\to\infty$.

\paragraph{NEW (combination with a KNOWN estimate).}
\begin{maincorollary}[A positive rate strictly below the Diophantine rate]\label{cor:bracket}
There is an absolute constant $C\ge4$ such that for $\lambda>1$ and irrational $\alpha$ with $0<x\le2/C$,
\[
 \frac1C\le\rho(\lambda,\alpha)\le1-x<1.
\]
Here $C$ can be any sufficiently enlarged constant in Liu--Shi's Theorem~1.3~\cite{LiuShi2019}. Its numerical value is not specified.
\end{maincorollary}

\paragraph{KNOWN (version~1.0.0 of this paper).}
We retain the results on the other side, including their proofs.
\begin{maintheorem}[Polynomial lower subsequences]\label{thm:polynomial}
Suppose $L/2<\beta<\infty$.
\begin{enumerate}
\item If $\beta<3L/2$, for every $\beta'\in(L/2,\beta)$ and $\epsilon>0$ there are infinitely many labels $m$ with
\[
 |G_m|\ge |m|^{-\frac{\beta'+L/2}{\beta'-L/2}-\epsilon}.
\]
\item If $3L/2\le\beta<\infty$, for every $\epsilon>0$ there are infinitely many labels $m$ with
\[
 |G_m|\ge |m|^{-\frac{\beta+L/2}{L}-\epsilon}.
\]
\end{enumerate}
In both cases $\rho=0$. The labels can be chosen as $m=-Dq_n\sgn\delta_n$ at the top of the spectrum.
\end{maintheorem}

\paragraph{KNOWN (version~1.0.0 of this paper).}
The scale estimate behind the preceding theorem permits growing $D$.
\begin{maintheorem}[Gaps at one Liouville scale]\label{thm:scale}
There are absolute constants $c_0,c_1>0$ with the following property. Fix $\xi\in(0,L]$. For all sufficiently large $n$ with $b_n\ge L/2+\xi$, every integer
\begin{equation}\label{eq:Dcap}
1\le D\le
\left\lfloor\min\left\{\frac{c_1}{\hbar_n},
                 \Lambda_ne^{-\xi q_n}\right\}\right\rfloor
\end{equation}
has a gap with
\begin{equation}\label{eq:width}
 N|_{G_m}=1-D|\delta_n|,\quad m=-Dq_n\sgn\delta_n,\qquad
 |G_m|\ge\frac{c_0\hbar_n}{P_n'(E_n)}.
\end{equation}
For each $\eta>0$ the width bound is at least
$C_\eta\exp(-(b_n+L/2+\eta)q_n)$ eventually, with $C_\eta>0$ independent of $b_n,D$.
Equivalently it is $\exp(-(b_n+L/2+o(1))q_n)$, uniformly in $b_n$.
The first cap is smaller when $b_n\le3L/2-\xi$ and $n$ is large.
\end{maintheorem}
The change at $3L/2$ records comparison precision, not a claimed spectral transition. The separate surplus argument in Section~\ref{sec:surplus} covers $\beta=\infty$.

\subsection{Relation to earlier work}
\paragraph{KNOWN.}
All noncritical allowed gaps are open by Avila--You--Zhou~\cite[Theorem~1.1]{AYZdry}. Our results concern widths. Liu--Shi~\cite[Theorem~1.3]{LiuShi2019} prove exponential gap upper bounds at sufficiently small dual coupling. Section~\ref{sec:below-literature} proves the label-preserving duality passage and quotes their Remark~1.4(4) exactly. Theorem~\ref{thm:below} is contrary to the expectation in that remark on the nonempty overlap $0<x\le2/C$. It does not contradict their theorem, and it does not prove $\rho=1-x$.

\paragraph{KNOWN IN A SPECIAL CASE (Diophantine frequencies).}
Ge--You--Zhou~\cite[Theorem~1.1]{GYZ} prove $\rho=1$ for Diophantine $\alpha$; earlier estimates are due to Leguil--You--Zhao--Zhou~\cite{LYZZ2017}. These statements do not settle every frequency with $\beta=0$. Avila--You--Zhou~\cite[Theorem~6.2]{AYZdry} also give rational-gap estimates in a different label range. At critical coupling, Krasovsky~\cite[Theorems~3--4]{Krasovsky} obtains polynomial inherited-gap bounds with approximation and parity hypotheses. Polynomial Liouville-gap phenomena themselves are therefore known.

\paragraph{KNOWN (the reduction mechanisms).}
Wilkinson~\cite{Wilkinson1987} develops projected Hamiltonians near rational magnetic flux. Helffer--Sj\"ostrand~\cite[Section~5.1]{HSII} give a rigorous one-step reduction with translated reference wells; their iterated results have additional arithmetic hypotheses. The orbital displacement is the Rammal--Wilkinson correction~\cite{RammalBellissard,BellissardBarelli}. Gram orthogonalization and interaction matrices are standard multiple-well tools~\cite[Sections~2.1--2.2 and 3.2]{HSwells}. The IDS limit is classical~\cite{AvronSimon}.

\paragraph{NEW (the quantitative combination).}
Below the threshold we retain the full energy-dependent self-energy. A specified continuation removes remote physical wells and produces one covariant family. A relative walk estimate resolves gaps of order $\kappa^{1-D}$ at each fixed $D$. A common spectral margin and total Schur error then give an exact integer count in every long box. Above the threshold the retained proof uses common Mathieu windows with growing $D$. Neither argument claims priority for its individual mechanisms.

\subsection{Reading the proof}
\paragraph{NEW / KNOWN (proof organization).}
Section~\ref{sec:below} gives the fixed-$D$ reduction and the Liu--Shi comparison. The earlier above-threshold conclusions and surplus proof remain in Sections~\ref{sec:consequences}--\ref{sec:surplus}. The phase diagram separates the new gap bound from known spectral and transport statements. Appendices~\ref{app:complex}--\ref{app:limits} retain the earlier machinery. Appendices~\ref{app:below-resolvent}--\ref{app:below-assembly} prove the new continuation, self-energy estimates, covariant model theorem and IDS assembly. Appendix~\ref{app:companions} states every companion import with its source.
\label{end:intro}
